\Chapter{63}

\Verse{1}
{
    \zA{话说宝玉回至房中洗手，因和袭人商议：“晚间吃酒，大家取乐，不可拘泥。}
}
{
    \eA{[English source not present in the checked-in Bencraft/Joly files.]}
}
{
}

\Verse{2}
{
    \zA{如今吃什么好？}
}
{
    \eA{[English source not present in the checked-in Bencraft/Joly files.]}
}
{
}

\Verse{3}
{
    \zA{早说给他们备办去。”}
}
{
    \eA{[English source not present in the checked-in Bencraft/Joly files.]}
}
{
}

\Verse{4}
{
    \zA{袭人笑道：“你放心，我和晴雯、麝月、秋 纹四个人，每人五钱银子，共是二两；芳官、碧痕、春燕、四儿四个人，每人三钱
        银子，他们告假的不算：共是三两二钱银子，早已交给了柳嫂子，预备四十碟果子。}
}
{
    \eA{[English source not present in the checked-in Bencraft/Joly files.]}
}
{
}

\Verse{5}
{
    \zA{我和平儿说了，已经抬了一罐好绍兴酒藏在那边了。}
}
{
    \eA{[English source not present in the checked-in Bencraft/Joly files.]}
}
{
}

\Verse{6}
{
    \zA{我们八个人单替你做生日。”}
}
{
    \eA{[English source not present in the checked-in Bencraft/Joly files.]}
}
{
}

\Verse{7}
{
    \zA{宝玉听了，喜的忙说：“他们是那里的钱？}
}
{
    \eA{[English source not present in the checked-in Bencraft/Joly files.]}
}
{
}

\Verse{8}
{
    \zA{不该叫他们出才是。”}
}
{
    \eA{[English source not present in the checked-in Bencraft/Joly files.]}
}
{
}

\Verse{9}
{
    \zA{晴雯道：“他们 没钱，难道我们是有钱的？}
}
{
    \eA{[English source not present in the checked-in Bencraft/Joly files.]}
}
{
}

\Verse{10}
{
    \zA{这原是各人的心。}
}
{
    \eA{[English source not present in the checked-in Bencraft/Joly files.]}
}
{
}

\Verse{11}
{
    \zA{那怕他偷的呢，只管领他的情就是了。”}
}
{
    \eA{[English source not present in the checked-in Bencraft/Joly files.]}
}
{
}

\Verse{12}
{
    \zA{宝玉听了，笑说：“你说的是。”}
}
{
    \eA{[English source not present in the checked-in Bencraft/Joly files.]}
}
{
}

\Verse{13}
{
    \zA{袭人笑道：“你这个人，一天不捱他两句硬话村 你，你再过不去。”}
}
{
    \eA{[English source not present in the checked-in Bencraft/Joly files.]}
}
{
}

\Verse{14}
{
    \zA{晴雯笑道：“你如今也学坏了，专会调三窝四。”}
}
{
    \eA{[English source not present in the checked-in Bencraft/Joly files.]}
}
{
}

\Verse{15}
{
    \zA{说着，大家 都笑了。}
}
{
    \eA{[English source not present in the checked-in Bencraft/Joly files.]}
}
{
}

\Verse{16}
{
    \zA{宝玉说：“关了院门罢。”}
}
{
    \eA{[English source not present in the checked-in Bencraft/Joly files.]}
}
{
}

\Verse{17}
{
    \zA{袭人笑道：“怪不得人说你是‘无事忙’!这 会子关了门，人倒疑惑起来，索性再等一等。”}
}
{
    \eA{[English source not present in the checked-in Bencraft/Joly files.]}
}
{
}

\Verse{18}
{
    \zA{宝玉点头，因说：“我出去走走。}
}
{
    \eA{[English source not present in the checked-in Bencraft/Joly files.]}
}
{
}

\Verse{19}
{
    \zA{四儿舀水去，春燕一个跟我来罢。”}
}
{
    \eA{[English source not present in the checked-in Bencraft/Joly files.]}
}
{
}

\Verse{20}
{
    \zA{说着，走至外边，因见无人，便问五儿之事。}
}
{
    \eA{[English source not present in the checked-in Bencraft/Joly files.]}
}
{
}

\Verse{21}
{
    \zA{春燕道：“我才告诉了柳嫂子，他倒很喜欢。}
}
{
    \eA{[English source not present in the checked-in Bencraft/Joly files.]}
}
{
}

\Verse{22}
{
    \zA{只是五儿那一夜受了委屈烦恼，回去 又气病了，那里来得？}
}
{
    \eA{[English source not present in the checked-in Bencraft/Joly files.]}
}
{
}

\Verse{23}
{
    \zA{只等好了罢。”}
}
{
    \eA{[English source not present in the checked-in Bencraft/Joly files.]}
}
{
}

\Verse{24}
{
    \zA{宝玉听了，未免后悔长叹，因又问：“这事 袭人知道不知道？”}
}
{
    \eA{[English source not present in the checked-in Bencraft/Joly files.]}
}
{
}

\Verse{25}
{
    \zA{春燕道：“我没告诉，不知芳官可说了没有。”}
}
{
    \eA{[English source not present in the checked-in Bencraft/Joly files.]}
}
{
}

\Verse{26}
{
    \zA{宝玉道：“我 却没告诉过他。}
}
{
    \eA{[English source not present in the checked-in Bencraft/Joly files.]}
}
{
}

\Verse{27}
{
    \zA{也罢，等我告诉他就是了。”}
}
{
    \eA{[English source not present in the checked-in Bencraft/Joly files.]}
}
{
}

\Verse{28}
{
    \zA{说毕，复走进来，故意洗手。}
}
{
    \eA{[English source not present in the checked-in Bencraft/Joly files.]}
}
{
}

\Verse{29}
{
    \zA{已是掌灯时分，听得院门前有一群人进来。}
}
{
    \eA{[English source not present in the checked-in Bencraft/Joly files.]}
}
{
}

\Verse{30}
{
    \zA{大家隔窗悄视，果见林之孝家的和 几个管事的女人走来，前头一人提着大灯笼。}
}
{
    \eA{[English source not present in the checked-in Bencraft/Joly files.]}
}
{
}

\Verse{31}
{
    \zA{晴雯悄笑道：“他们查上夜的人来了。}
}
{
    \eA{[English source not present in the checked-in Bencraft/Joly files.]}
}
{
}

\Verse{32}
{
    \zA{这一出去，咱们就好关门了。”}
}
{
    \eA{[English source not present in the checked-in Bencraft/Joly files.]}
}
{
}

\Verse{33}
{
    \zA{只见怡红院凡上夜的人，都迎出去了。}
}
{
    \eA{[English source not present in the checked-in Bencraft/Joly files.]}
}
{
}

\Verse{34}
{
    \zA{林之孝家的 看了不少，又吩咐：“别耍钱吃酒，放倒头睡到大天亮。}
}
{
    \eA{[English source not present in the checked-in Bencraft/Joly files.]}
}
{
}

\Verse{35}
{
    \zA{我听见是不依的。”}
}
{
    \eA{[English source not present in the checked-in Bencraft/Joly files.]}
}
{
}

\Verse{36}
{
    \zA{众人 都笑说：“那里有这么大胆子的人。”}
}
{
    \eA{[English source not present in the checked-in Bencraft/Joly files.]}
}
{
}

\Verse{37}
{
    \zA{林之孝家的又问：“宝二爷睡下了没有？”}
}
{
    \eA{[English source not present in the checked-in Bencraft/Joly files.]}
}
{
}

\Verse{38}
{
    \zA{众人都回：“不知道。”}
}
{
    \eA{[English source not present in the checked-in Bencraft/Joly files.]}
}
{
}

\Verse{39}
{
    \zA{袭人忙推宝玉。}
}
{
    \eA{[English source not present in the checked-in Bencraft/Joly files.]}
}
{
}

\Verse{40}
{
    \zA{宝玉？}
}
{
    \eA{[English source not present in the checked-in Bencraft/Joly files.]}
}
{
}

\Verse{41}
{
    \zA{了鞋，便迎出来，笑道：“我还没 睡呢。}
}
{
    \eA{[English source not present in the checked-in Bencraft/Joly files.]}
}
{
}

\Verse{42}
{
    \zA{妈妈进来歇歇。”}
}
{
    \eA{[English source not present in the checked-in Bencraft/Joly files.]}
}
{
}

\Verse{43}
{
    \zA{又叫：“袭人，倒茶来。”}
}
{
    \eA{[English source not present in the checked-in Bencraft/Joly files.]}
}
{
}

\Verse{44}
{
    \zA{林之孝家的忙进来，笑说：“还 没睡呢？}
}
{
    \eA{[English source not present in the checked-in Bencraft/Joly files.]}
}
{
}

\Verse{45}
{
    \zA{如今天长夜短，该早些睡了，明日方起的早。}
}
{
    \eA{[English source not present in the checked-in Bencraft/Joly files.]}
}
{
}

\Verse{46}
{
    \zA{不然，到了明日起迟了，人 家笑话，不是个读书上学的公子了，倒像那起挑脚汉了。”}
}
{
    \eA{[English source not present in the checked-in Bencraft/Joly files.]}
}
{
}

\Verse{47}
{
    \zA{说毕，又笑。}
}
{
    \eA{[English source not present in the checked-in Bencraft/Joly files.]}
}
{
}

\Verse{48}
{
    \zA{宝玉忙笑 道：“妈妈说的是。}
}
{
    \eA{[English source not present in the checked-in Bencraft/Joly files.]}
}
{
}

\Verse{49}
{
    \zA{我每日都睡的早，妈妈每日进来，可都是我不知道的，已经睡 了。}
}
{
    \eA{[English source not present in the checked-in Bencraft/Joly files.]}
}
{
}

\Verse{50}
{
    \zA{今日因吃了面，怕停食，所以多玩一回。”}
}
{
    \eA{[English source not present in the checked-in Bencraft/Joly files.]}
}
{
}

\Verse{51}
{
    \zA{林之孝家的又向袭人等笑说：“该 ？}
}
{
    \eA{[English source not present in the checked-in Bencraft/Joly files.]}
}
{
}

\Verse{52}
{
    \zA{些普洱茶喝。”}
}
{
    \eA{[English source not present in the checked-in Bencraft/Joly files.]}
}
{
}

\Verse{53}
{
    \zA{袭人晴雯二人忙说：“？}
}
{
    \eA{[English source not present in the checked-in Bencraft/Joly files.]}
}
{
}

\Verse{54}
{
    \zA{了一茶缸子女儿茶，已经喝过两碗了。}
}
{
    \eA{[English source not present in the checked-in Bencraft/Joly files.]}
}
{
}

\Verse{55}
{
    \zA{大娘也尝一碗，都是现成的。”}
}
{
    \eA{[English source not present in the checked-in Bencraft/Joly files.]}
}
{
}

\Verse{56}
{
    \zA{说着，晴雯便倒了来。}
}
{
    \eA{[English source not present in the checked-in Bencraft/Joly files.]}
}
{
}

\Verse{57}
{
    \zA{林家的站起接了，又笑道： “这些时，我听见二爷嘴里都换了字眼，赶着这几位大姑娘们竟叫起名字来。}
}
{
    \eA{[English source not present in the checked-in Bencraft/Joly files.]}
}
{
}

\Verse{58}
{
    \zA{虽然 在这屋里，到底是老太太、太太的人，还该嘴里尊重些才是。}
}
{
    \eA{[English source not present in the checked-in Bencraft/Joly files.]}
}
{
}

\Verse{59}
{
    \zA{若一时半刻偶然叫一 声使得；若只管顺口叫起来，怕以后兄弟侄儿照样，就惹人笑话这家子的人眼里没 有长辈了。”}
}
{
    \eA{[English source not present in the checked-in Bencraft/Joly files.]}
}
{
}

\Verse{60}
{
    \zA{宝玉笑道：“妈妈说的是。}
}
{
    \eA{[English source not present in the checked-in Bencraft/Joly files.]}
}
{
}

\Verse{61}
{
    \zA{我不过是一时半刻偶然叫一句是有的。”}
}
{
    \eA{[English source not present in the checked-in Bencraft/Joly files.]}
}
{
}

\Verse{62}
{
    \zA{袭人晴雯都笑说：“这可别委屈了他，直到如今，他可‘姐姐’没离了嘴。}
}
{
    \eA{[English source not present in the checked-in Bencraft/Joly files.]}
}
{
}

\Verse{63}
{
    \zA{不过玩 的时候叫一声半声名字，若当着人，却是和先一样。”}
}
{
    \eA{[English source not present in the checked-in Bencraft/Joly files.]}
}
{
}

\Verse{64}
{
    \zA{林之孝家的笑道：“这才好 呢，这才是读书知礼的。}
}
{
    \eA{[English source not present in the checked-in Bencraft/Joly files.]}
}
{
}

\Verse{65}
{
    \zA{越自己谦逊，越尊重。}
}
{
    \eA{[English source not present in the checked-in Bencraft/Joly files.]}
}
{
}

\Verse{66}
{
    \zA{别说是三五代的陈人、现从老太太、 太太屋里拨过来的，就是老太太、太太屋里的猫儿狗儿，轻易也伤不得他。}
}
{
    \eA{[English source not present in the checked-in Bencraft/Joly files.]}
}
{
}

\Verse{67}
{
    \zA{这才是 受过调教的公子行事。”}
}
{
    \eA{[English source not present in the checked-in Bencraft/Joly files.]}
}
{
}

\Verse{68}
{
    \zA{说毕，吃了茶，便说：“请安歇罢，我们走了。”}
}
{
    \eA{[English source not present in the checked-in Bencraft/Joly files.]}
}
{
}

\Verse{69}
{
    \zA{宝玉还 说：“再歇歇。”}
}
{
    \eA{[English source not present in the checked-in Bencraft/Joly files.]}
}
{
}

\Verse{70}
{
    \zA{那林之孝家的已带了众人又查别处去了。}
}
{
    \eA{[English source not present in the checked-in Bencraft/Joly files.]}
}
{
}

\Verse{71}
{
    \zA{这里晴雯等忙命关了门，进来笑说：“这位奶奶那里吃了一杯来了？}
}
{
    \eA{[English source not present in the checked-in Bencraft/Joly files.]}
}
{
}

\Verse{72}
{
    \zA{唠三叨四 的，又排场了我们一顿去了。”}
}
{
    \eA{[English source not present in the checked-in Bencraft/Joly files.]}
}
{
}

\Verse{73}
{
    \zA{麝月笑道：“他也不是好意的？}
}
{
    \eA{[English source not present in the checked-in Bencraft/Joly files.]}
}
{
}

\Verse{74}
{
    \zA{少不得也要常提着 些儿，也堤防着，怕走了大褶儿的意思。”}
}
{
    \eA{[English source not present in the checked-in Bencraft/Joly files.]}
}
{
}

\Verse{75}
{
    \zA{说着，一面摆上酒果。}
}
{
    \eA{[English source not present in the checked-in Bencraft/Joly files.]}
}
{
}

\Verse{76}
{
    \zA{袭人道：“不用 高桌，咱们把那张花梨圆炕桌子放在炕上坐，又宽绰，又便宜。”}
}
{
    \eA{[English source not present in the checked-in Bencraft/Joly files.]}
}
{
}

\Verse{77}
{
    \zA{说着，大家果然 抬来。}
}
{
    \eA{[English source not present in the checked-in Bencraft/Joly files.]}
}
{
}

\Verse{78}
{
    \zA{麝月和四儿那边去搬果子，用两个大茶盘，做四五次方搬运了来。}
}
{
    \eA{[English source not present in the checked-in Bencraft/Joly files.]}
}
{
}

\Verse{79}
{
    \zA{两个老婆 子蹲在外面火盆上筛酒。}
}
{
    \eA{[English source not present in the checked-in Bencraft/Joly files.]}
}
{
}

\Verse{80}
{
    \zA{宝玉说：“天热，咱们都脱了大衣裳才好。”}
}
{
    \eA{[English source not present in the checked-in Bencraft/Joly files.]}
}
{
}

\Verse{81}
{
    \zA{众人笑道： “你要脱，你脱，我们还要轮流安席呢。”}
}
{
    \eA{[English source not present in the checked-in Bencraft/Joly files.]}
}
{
}

\Verse{82}
{
    \zA{宝玉笑道：“这一安席，就要到五更天 了。}
}
{
    \eA{[English source not present in the checked-in Bencraft/Joly files.]}
}
{
}

\Verse{83}
{
    \zA{知道我最怕这些俗套，在外人跟前，不得已的。}
}
{
    \eA{[English source not present in the checked-in Bencraft/Joly files.]}
}
{
}

\Verse{84}
{
    \zA{这会子还怄我，就不好了。”}
}
{
    \eA{[English source not present in the checked-in Bencraft/Joly files.]}
}
{
}

\Verse{85}
{
    \zA{众人听了，都说：“依你。”}
}
{
    \eA{[English source not present in the checked-in Bencraft/Joly files.]}
}
{
}

\Verse{86}
{
    \zA{于是先不上坐，且忙着卸妆宽衣。}
}
{
    \eA{[English source not present in the checked-in Bencraft/Joly files.]}
}
{
}

\Verse{87}
{
    \zA{一时将正妆卸去，头上只随便挽着？}
}
{
    \eA{[English source not present in the checked-in Bencraft/Joly files.]}
}
{
}

\Verse{88}
{
    \zA{儿，身 上皆是紧身袄儿。}
}
{
    \eA{[English source not present in the checked-in Bencraft/Joly files.]}
}
{
}

\Verse{89}
{
    \zA{宝玉只穿着大红绵纱小袄儿，下面绿绫弹墨夹裤，散着裤脚，系 着一条汗巾，靠着一个各色玫瑰芍药花瓣装的玉色夹纱新枕头，和芳官两个先？}
}
{
    \eA{[English source not present in the checked-in Bencraft/Joly files.]}
}
{
}

\Verse{90}
{
    \zA{拳。}
}
{
    \eA{[English source not present in the checked-in Bencraft/Joly files.]}
}
{
}

\Verse{91}
{
    \zA{当时芳官满口嚷热，只穿着一件玉色红青驼绒三色缎子拼的水田小夹袄，束着 一条柳绿汗巾，底下是水红洒花夹裤，也散着裤腿。}
}
{
    \eA{[English source not present in the checked-in Bencraft/Joly files.]}
}
{
}

\Verse{92}
{
    \zA{头上齐额编着一圈小辫，总归 至顶心，结一根粗辫，拖在脑后，右耳根内只塞着米粒大小的一个小玉塞子，左耳
        上单一个白果大小的硬红镶金大坠子，越显得面如满月犹白，眼似秋水还清。}
}
{
    \eA{[English source not present in the checked-in Bencraft/Joly files.]}
}
{
}

\Verse{93}
{
    \zA{引得 众人笑说：“他两个倒像一对双生的弟兄。”}
}
{
    \eA{[English source not present in the checked-in Bencraft/Joly files.]}
}
{
}

\Verse{94}
{
    \zA{袭人等一一斟上酒来，说：“且等一 等再？}
}
{
    \eA{[English source not present in the checked-in Bencraft/Joly files.]}
}
{
}

\Verse{95}
{
    \zA{拳。}
}
{
    \eA{[English source not present in the checked-in Bencraft/Joly files.]}
}
{
}

\Verse{96}
{
    \zA{虽不安席，在我们每人手里吃一口罢了。”}
}
{
    \eA{[English source not present in the checked-in Bencraft/Joly files.]}
}
{
}

\Verse{97}
{
    \zA{于是袭人为先，端在唇上吃 了一口，其馀依次下去，一一吃过，大家方团圆坐了。}
}
{
    \eA{[English source not present in the checked-in Bencraft/Joly files.]}
}
{
}

\Verse{98}
{
    \zA{春燕四儿因炕沿坐不下，便 端了两个绒套绣墩近炕沿放下。}
}
{
    \eA{[English source not present in the checked-in Bencraft/Joly files.]}
}
{
}

\Verse{99}
{
    \zA{那四十个碟子，皆是一色白彩定窑的，不过小茶碟 大，里面自是山南海北干鲜水陆的酒馔果菜。}
}
{
    \eA{[English source not present in the checked-in Bencraft/Joly files.]}
}
{
}

\Verse{100}
{
    \zA{宝玉因说：“咱们也该行个令才好。”}
}
{
    \eA{[English source not present in the checked-in Bencraft/Joly files.]}
}
{
}

\Verse{101}
{
    \zA{袭人道：“斯文些才好，别大呼小叫， 叫人听见。}
}
{
    \eA{[English source not present in the checked-in Bencraft/Joly files.]}
}
{
}

\Verse{102}
{
    \zA{二则我们不识字，可不要那些文的。”}
}
{
    \eA{[English source not present in the checked-in Bencraft/Joly files.]}
}
{
}

\Verse{103}
{
    \zA{麝月笑道：“拿骰子咱们抢红罢。”}
}
{
    \eA{[English source not present in the checked-in Bencraft/Joly files.]}
}
{
}

\Verse{104}
{
    \zA{宝玉道：“没趣，不好。}
}
{
    \eA{[English source not present in the checked-in Bencraft/Joly files.]}
}
{
}

\Verse{105}
{
    \zA{咱们占花名儿好。”}
}
{
    \eA{[English source not present in the checked-in Bencraft/Joly files.]}
}
{
}

\Verse{106}
{
    \zA{晴雯笑道：“正是，早已想弄这个玩 意儿。”}
}
{
    \eA{[English source not present in the checked-in Bencraft/Joly files.]}
}
{
}

\Verse{107}
{
    \zA{袭人道：“这个玩意虽好，人少了没趣。”}
}
{
    \eA{[English source not present in the checked-in Bencraft/Joly files.]}
}
{
}

\Verse{108}
{
    \zA{春燕笑道：“依我说，咱们竟 悄悄的把宝姑娘、云姑娘、林姑娘请了来，玩一会子，到二更天再睡不迟。”}
}
{
    \eA{[English source not present in the checked-in Bencraft/Joly files.]}
}
{
}

\Verse{109}
{
    \zA{袭人 道：“又开门合户的闹，倘或遇见巡夜的问？”}
}
{
    \eA{[English source not present in the checked-in Bencraft/Joly files.]}
}
{
}

\Verse{110}
{
    \zA{宝玉道：“怕什么!咱们三姑娘也 吃酒，再请他一声才好。}
}
{
    \eA{[English source not present in the checked-in Bencraft/Joly files.]}
}
{
}

\Verse{111}
{
    \zA{还有琴姑娘。”}
}
{
    \eA{[English source not present in the checked-in Bencraft/Joly files.]}
}
{
}

\Verse{112}
{
    \zA{众人都道：“琴姑娘罢了，他在大奶奶屋 里，叨登的大发了。”}
}
{
    \eA{[English source not present in the checked-in Bencraft/Joly files.]}
}
{
}

\Verse{113}
{
    \zA{宝玉道：“怕什么，你们就快请去。”}
}
{
    \eA{[English source not present in the checked-in Bencraft/Joly files.]}
}
{
}

\Verse{114}
{
    \zA{春燕四儿都巴不得一 声，二人忙命开门，各带小丫头分头去请。}
}
{
    \eA{[English source not present in the checked-in Bencraft/Joly files.]}
}
{
}

\Verse{115}
{
    \zA{晴雯、麝月、袭人三人又说：“他两个去请，只怕不肯来，须得我们去请，死 活拉了来。”}
}
{
    \eA{[English source not present in the checked-in Bencraft/Joly files.]}
}
{
}

\Verse{116}
{
    \zA{于是袭人晴雯忙又命老婆子打个灯笼，二人又去。}
}
{
    \eA{[English source not present in the checked-in Bencraft/Joly files.]}
}
{
}

\Verse{117}
{
    \zA{果然宝钗说夜深了， 黛玉说身上不好。}
}
{
    \eA{[English source not present in the checked-in Bencraft/Joly files.]}
}
{
}

\Verse{118}
{
    \zA{他二人再三央求：“好歹给我们一点体面，略坐坐再来。”}
}
{
    \eA{[English source not present in the checked-in Bencraft/Joly files.]}
}
{
}

\Verse{119}
{
    \zA{众人 听了，却也欢喜。}
}
{
    \eA{[English source not present in the checked-in Bencraft/Joly files.]}
}
{
}

\Verse{120}
{
    \zA{因想不请李纨，倘或被他知道了倒不好，便命翠墨同春燕也再三 的请了李纨和宝琴二人，会齐先后都到了怡红院中。}
}
{
    \eA{[English source not present in the checked-in Bencraft/Joly files.]}
}
{
}

\Verse{121}
{
    \zA{袭人又死活拉了香菱来。}
}
{
    \eA{[English source not present in the checked-in Bencraft/Joly files.]}
}
{
}

\Verse{122}
{
    \zA{炕上 又并了一张桌子，方坐开了。}
}
{
    \eA{[English source not present in the checked-in Bencraft/Joly files.]}
}
{
}

\Verse{123}
{
    \zA{宝玉忙说：“林妹妹怕冷，过这边靠板壁坐。”}
}
{
    \eA{[English source not present in the checked-in Bencraft/Joly files.]}
}
{
}

\Verse{124}
{
    \zA{又拿 了个靠背垫着些。}
}
{
    \eA{[English source not present in the checked-in Bencraft/Joly files.]}
}
{
}

\Verse{125}
{
    \zA{袭人等都端了椅子在炕沿下陪着。}
}
{
    \eA{[English source not present in the checked-in Bencraft/Joly files.]}
}
{
}

\Verse{126}
{
    \zA{黛玉却离桌远远的靠着靠背， 因笑向宝钗、李纨、探春等道：“你们日日说人家夜饮聚赌，今日我们自己也如此。}
}
{
    \eA{[English source not present in the checked-in Bencraft/Joly files.]}
}
{
}

\Verse{127}
{
    \zA{以后怎么说人？”}
}
{
    \eA{[English source not present in the checked-in Bencraft/Joly files.]}
}
{
}

\Verse{128}
{
    \zA{李纨笑道：“有何妨碍？}
}
{
    \eA{[English source not present in the checked-in Bencraft/Joly files.]}
}
{
}

\Verse{129}
{
    \zA{一年之中不过生日节间如此，并没夜夜 如此，这倒也不怕。”}
}
{
    \eA{[English source not present in the checked-in Bencraft/Joly files.]}
}
{
}

\Verse{130}
{
    \zA{说着，晴雯拿了一个竹雕的签筒来，里面装着象牙花名签子，摇了一摇，放在 当中。}
}
{
    \eA{[English source not present in the checked-in Bencraft/Joly files.]}
}
{
}

\Verse{131}
{
    \zA{又取过骰子来，盛在盒内，摇了一摇，揭开一看，里面是六点，数至宝钗。}
}
{
    \eA{[English source not present in the checked-in Bencraft/Joly files.]}
}
{
}

\Verse{132}
{
    \zA{宝钗便笑道：“我先抓，不知抓出个什么来。”}
}
{
    \eA{[English source not present in the checked-in Bencraft/Joly files.]}
}
{
}

\Verse{133}
{
    \zA{说着将筒摇了一摇，伸手掣出一签。}
}
{
    \eA{[English source not present in the checked-in Bencraft/Joly files.]}
}
{
}

\Verse{134}
{
    \zA{大家一看，只见签上画着一枝牡丹，题着“艳冠群芳”四字。}
}
{
    \eA{[English source not present in the checked-in Bencraft/Joly files.]}
}
{
}

\Verse{135}
{
    \zA{下面又有镌的小字， 一句唐诗，道是： 任是无情也动人。}
}
{
    \eA{[English source not present in the checked-in Bencraft/Joly files.]}
}
{
}

\Verse{136}
{
    \zA{又注着：“在席共贺一杯。}
}
{
    \eA{[English source not present in the checked-in Bencraft/Joly files.]}
}
{
}

\Verse{137}
{
    \zA{此为群芳之冠，随意命人，不拘诗词雅谑，或新曲一支 为贺。”}
}
{
    \eA{[English source not present in the checked-in Bencraft/Joly files.]}
}
{
}

\Verse{138}
{
    \zA{众人都笑说：“巧得很!你也原配牡丹花。”}
}
{
    \eA{[English source not present in the checked-in Bencraft/Joly files.]}
}
{
}

\Verse{139}
{
    \zA{说着大家共贺了一杯。}
}
{
    \eA{[English source not present in the checked-in Bencraft/Joly files.]}
}
{
}

\Verse{140}
{
    \zA{宝钗 吃过，便笑说：“芳官唱一只我们听罢。”}
}
{
    \eA{[English source not present in the checked-in Bencraft/Joly files.]}
}
{
}

\Verse{141}
{
    \zA{芳官道：“既这样，大家吃了门杯好听。”}
}
{
    \eA{[English source not present in the checked-in Bencraft/Joly files.]}
}
{
}

\Verse{142}
{
    \zA{于是大家吃酒，芳官便唱：“寿筵开处风光好……”众人都道：“快打回去!这会 子很不用你来上寿。}
}
{
    \eA{[English source not present in the checked-in Bencraft/Joly files.]}
}
{
}

\Verse{143}
{
    \zA{拣你极好的唱来。”}
}
{
    \eA{[English source not present in the checked-in Bencraft/Joly files.]}
}
{
}

\Verse{144}
{
    \zA{芳官只得细细的唱了一只《赏花时》“翠 凤翎毛扎帚？}
}
{
    \eA{[English source not present in the checked-in Bencraft/Joly files.]}
}
{
}

\Verse{145}
{
    \zA{，闲踏天门扫落花……”才罢。}
}
{
    \eA{[English source not present in the checked-in Bencraft/Joly files.]}
}
{
}

\Verse{146}
{
    \zA{宝玉却只管拿着那签，口内颠来倒去 念“任是无情也动人”，听了这曲子，眼看着芳官不语。}
}
{
    \eA{[English source not present in the checked-in Bencraft/Joly files.]}
}
{
}

\Verse{147}
{
    \zA{湘云忙一手夺了，撂与宝 钗。}
}
{
    \eA{[English source not present in the checked-in Bencraft/Joly files.]}
}
{
}

\Verse{148}
{
    \zA{宝钗又掷了一个十六点，数到探春。}
}
{
    \eA{[English source not present in the checked-in Bencraft/Joly files.]}
}
{
}

\Verse{149}
{
    \zA{探春笑道：“还不知得个什么。”}
}
{
    \eA{[English source not present in the checked-in Bencraft/Joly files.]}
}
{
}

\Verse{150}
{
    \zA{伸手掣 了一根出来，自己一瞧，便撂在桌上，红了脸笑道：“很不该行这个令!这原是外 头男人们行的令，许多混帐话在上头。”}
}
{
    \eA{[English source not present in the checked-in Bencraft/Joly files.]}
}
{
}

\Verse{151}
{
    \zA{众人不解，袭人等忙拾起来。}
}
{
    \eA{[English source not present in the checked-in Bencraft/Joly files.]}
}
{
}

\Verse{152}
{
    \zA{众人看时， 上面一枝杏花，那红字写着“瑶池仙品”四字，诗云： 日边红杏倚云栽。}
}
{
    \eA{[English source not present in the checked-in Bencraft/Joly files.]}
}
{
}

\Verse{153}
{
    \zA{注云：“得此签者，必得贵婿，大家恭贺一杯，再同饮一杯。”}
}
{
    \eA{[English source not present in the checked-in Bencraft/Joly files.]}
}
{
}

\Verse{154}
{
    \zA{众人笑说道：“我 们说是什么呢，这签原是闺阁中取笑的，除了这两三根有这话的，并无杂话。}
}
{
    \eA{[English source not present in the checked-in Bencraft/Joly files.]}
}
{
}

\Verse{155}
{
    \zA{这有 何妨？}
}
{
    \eA{[English source not present in the checked-in Bencraft/Joly files.]}
}
{
}

\Verse{156}
{
    \zA{我们家已有了王妃，难道你也是王妃不成？}
}
{
    \eA{[English source not present in the checked-in Bencraft/Joly files.]}
}
{
}

\Verse{157}
{
    \zA{大喜，大喜！”}
}
{
    \eA{[English source not present in the checked-in Bencraft/Joly files.]}
}
{
}

\Verse{158}
{
    \zA{说着大家来敬探春。}
}
{
    \eA{[English source not present in the checked-in Bencraft/Joly files.]}
}
{
}

\Verse{159}
{
    \zA{探春那里肯饮，却被湘云、香菱、李纨等三四个人，强死强活，灌了一钟才罢。}
}
{
    \eA{[English source not present in the checked-in Bencraft/Joly files.]}
}
{
}

\Verse{160}
{
    \zA{探春只叫：“蠲了这个，再行别的。”}
}
{
    \eA{[English source not present in the checked-in Bencraft/Joly files.]}
}
{
}

\Verse{161}
{
    \zA{众人断不肯依。}
}
{
    \eA{[English source not present in the checked-in Bencraft/Joly files.]}
}
{
}

\Verse{162}
{
    \zA{湘云拿着他的手，强掷 了个十九点出来，便该李氏掣。}
}
{
    \eA{[English source not present in the checked-in Bencraft/Joly files.]}
}
{
}

\Verse{163}
{
    \zA{李氏摇了一摇，掣出一根来一看，笑道：“好极! 你们瞧瞧这行子，竟有些意思。”}
}
{
    \eA{[English source not present in the checked-in Bencraft/Joly files.]}
}
{
}

\Verse{164}
{
    \zA{众人瞧那签上，画着一枝老梅，写着“霜晓寒姿” 四字，那一面旧诗是： 竹篱茅舍自甘心。}
}
{
    \eA{[English source not present in the checked-in Bencraft/Joly files.]}
}
{
}

\Verse{165}
{
    \zA{注云：“自饮一杯，下家掷骰。”}
}
{
    \eA{[English source not present in the checked-in Bencraft/Joly files.]}
}
{
}

\Verse{166}
{
    \zA{李纨笑道：“真有趣，你们掷去罢，我只自吃一 杯，不问你们的废兴。”}
}
{
    \eA{[English source not present in the checked-in Bencraft/Joly files.]}
}
{
}

\Verse{167}
{
    \zA{说着便吃酒，将骰过给黛玉。}
}
{
    \eA{[English source not present in the checked-in Bencraft/Joly files.]}
}
{
}

\Verse{168}
{
    \zA{黛玉一掷是十八点，便该湘云掣。}
}
{
    \eA{[English source not present in the checked-in Bencraft/Joly files.]}
}
{
}

\Verse{169}
{
    \zA{湘云笑着，揎拳掳袖的，伸手掣了一根出来。}
}
{
    \eA{[English source not present in the checked-in Bencraft/Joly files.]}
}
{
}

\Verse{170}
{
    \zA{大家看时，一面画着一枝海棠，题着“香梦？}
}
{
    \eA{[English source not present in the checked-in Bencraft/Joly files.]}
}
{
}

\Verse{171}
{
    \zA{酣”四字。}
}
{
    \eA{[English source not present in the checked-in Bencraft/Joly files.]}
}
{
}

\Verse{172}
{
    \zA{那面诗道是： 只恐夜深花睡去。}
}
{
    \eA{[English source not present in the checked-in Bencraft/Joly files.]}
}
{
}

\Verse{173}
{
    \zA{黛玉笑道：“‘夜深’二字改‘石凉’两个字倒好。”}
}
{
    \eA{[English source not present in the checked-in Bencraft/Joly files.]}
}
{
}

\Verse{174}
{
    \zA{众人知他打趣日间湘云醉眠 的事，都笑了。}
}
{
    \eA{[English source not present in the checked-in Bencraft/Joly files.]}
}
{
}

\Verse{175}
{
    \zA{湘云笑指那自行船给黛玉看，又说：“快坐上那船家去罢，别多说 了。”}
}
{
    \eA{[English source not present in the checked-in Bencraft/Joly files.]}
}
{
}

\Verse{176}
{
    \zA{众人都笑了。}
}
{
    \eA{[English source not present in the checked-in Bencraft/Joly files.]}
}
{
}

\Verse{177}
{
    \zA{因看注云：“既云香梦？}
}
{
    \eA{[English source not present in the checked-in Bencraft/Joly files.]}
}
{
}

\Verse{178}
{
    \zA{酣，掣此签者，不便饮酒，只令上下 两家各饮一杯。”}
}
{
    \eA{[English source not present in the checked-in Bencraft/Joly files.]}
}
{
}

\Verse{179}
{
    \zA{湘云拍手笑道：“阿弥陀佛，真真好签！”}
}
{
    \eA{[English source not present in the checked-in Bencraft/Joly files.]}
}
{
}

\Verse{180}
{
    \zA{恰好黛玉是上家，宝 玉是下家，二人斟了两杯，只得要饮。}
}
{
    \eA{[English source not present in the checked-in Bencraft/Joly files.]}
}
{
}

\Verse{181}
{
    \zA{宝玉先饮了半杯，瞅人不见，递与芳官。}
}
{
    \eA{[English source not present in the checked-in Bencraft/Joly files.]}
}
{
}

\Verse{182}
{
    \zA{芳 官即便端起来，一仰脖喝了。}
}
{
    \eA{[English source not present in the checked-in Bencraft/Joly files.]}
}
{
}

\Verse{183}
{
    \zA{黛玉只管和人说话，将酒全折在漱盂内了。}
}
{
    \eA{[English source not present in the checked-in Bencraft/Joly files.]}
}
{
}

\Verse{184}
{
    \zA{湘云便抓起骰子来，一掷个九点，数去该麝月。}
}
{
    \eA{[English source not present in the checked-in Bencraft/Joly files.]}
}
{
}

\Verse{185}
{
    \zA{麝月便掣了一根出来，大家看 时，上面是一枝荼？}
}
{
    \eA{[English source not present in the checked-in Bencraft/Joly files.]}
}
{
}

\Verse{186}
{
    \zA{花，题着“韶华胜极”四字，那边写着一句旧诗，道是： 开到荼？}
}
{
    \eA{[English source not present in the checked-in Bencraft/Joly files.]}
}
{
}

\Verse{187}
{
    \zA{花事了。}
}
{
    \eA{[English source not present in the checked-in Bencraft/Joly files.]}
}
{
}

\Verse{188}
{
    \zA{注云：“在席各饮三杯送春。”}
}
{
    \eA{[English source not present in the checked-in Bencraft/Joly files.]}
}
{
}

\Verse{189}
{
    \zA{麝月问：“怎么讲？”}
}
{
    \eA{[English source not present in the checked-in Bencraft/Joly files.]}
}
{
}

\Verse{190}
{
    \zA{宝玉皱皱眉儿，忙将签藏了， 说：“咱们且喝酒罢。”}
}
{
    \eA{[English source not present in the checked-in Bencraft/Joly files.]}
}
{
}

\Verse{191}
{
    \zA{说着，大家吃了三口，以充三杯之数。}
}
{
    \eA{[English source not present in the checked-in Bencraft/Joly files.]}
}
{
}

\Verse{192}
{
    \zA{麝月一掷个十点，该香菱。}
}
{
    \eA{[English source not present in the checked-in Bencraft/Joly files.]}
}
{
}

\Verse{193}
{
    \zA{香菱便掣了一根并蒂花，题着“联春绕瑞”，那面 写着一句旧诗，道是： 连理枝头花正开。}
}
{
    \eA{[English source not present in the checked-in Bencraft/Joly files.]}
}
{
}

\Verse{194}
{
    \zA{注云：“共贺掣者三杯，大家陪饮一杯。”}
}
{
    \eA{[English source not present in the checked-in Bencraft/Joly files.]}
}
{
}

\Verse{195}
{
    \zA{香菱便又掷了个六点，该黛玉。}
}
{
    \eA{[English source not present in the checked-in Bencraft/Joly files.]}
}
{
}

\Verse{196}
{
    \zA{黛玉默默的想道：“不知还有什么好的被我掣 着方好。”}
}
{
    \eA{[English source not present in the checked-in Bencraft/Joly files.]}
}
{
}

\Verse{197}
{
    \zA{一面伸手取了一根。}
}
{
    \eA{[English source not present in the checked-in Bencraft/Joly files.]}
}
{
}

\Verse{198}
{
    \zA{只见上面画着一枝芙蓉花，题着“风露清愁”四字， 那面一句旧诗，道是： 莫怨东风当自嗟。}
}
{
    \eA{[English source not present in the checked-in Bencraft/Joly files.]}
}
{
}

\Verse{199}
{
    \zA{注云：“自饮一杯，牡丹陪饮一杯。”}
}
{
    \eA{[English source not present in the checked-in Bencraft/Joly files.]}
}
{
}

\Verse{200}
{
    \zA{众人笑说：“这个好极，除了他，别人不配 做芙蓉。”}
}
{
    \eA{[English source not present in the checked-in Bencraft/Joly files.]}
}
{
}

\Verse{201}
{
    \zA{黛玉也自笑了。}
}
{
    \eA{[English source not present in the checked-in Bencraft/Joly files.]}
}
{
}

\Verse{202}
{
    \zA{于是饮了酒，便掷了个二十点，该着袭人。}
}
{
    \eA{[English source not present in the checked-in Bencraft/Joly files.]}
}
{
}

\Verse{203}
{
    \zA{袭人便伸手取了一枝出来，却是一 枝桃花，题着“武陵别景”四字，那一面写着旧诗，道是： 桃红又见一年春。}
}
{
    \eA{[English source not present in the checked-in Bencraft/Joly files.]}
}
{
}

\Verse{204}
{
    \zA{注云：“杏花陪一盏，坐中同庚者陪一盏，同姓者陪一盏。”}
}
{
    \eA{[English source not present in the checked-in Bencraft/Joly files.]}
}
{
}

\Verse{205}
{
    \zA{众人笑道：“这一回 热闹有趣。”}
}
{
    \eA{[English source not present in the checked-in Bencraft/Joly files.]}
}
{
}

\Verse{206}
{
    \zA{大家算来：香菱、晴雯、宝钗三人皆与他同庚，黛玉与他同辰，只无 同姓者。}
}
{
    \eA{[English source not present in the checked-in Bencraft/Joly files.]}
}
{
}

\Verse{207}
{
    \zA{芳官忙道：“我也姓花，我也陪他一钟。”}
}
{
    \eA{[English source not present in the checked-in Bencraft/Joly files.]}
}
{
}

\Verse{208}
{
    \zA{于是大家斟了酒。}
}
{
    \eA{[English source not present in the checked-in Bencraft/Joly files.]}
}
{
}

\Verse{209}
{
    \zA{黛玉因向探 春笑道：“命中该招贵婿的!你是杏花，快喝了，我们好喝。”}
}
{
    \eA{[English source not present in the checked-in Bencraft/Joly files.]}
}
{
}

\Verse{210}
{
    \zA{探春笑道：“这是 什么话？}
}
{
    \eA{[English source not present in the checked-in Bencraft/Joly files.]}
}
{
}

\Verse{211}
{
    \zA{大嫂子顺手给他一巴掌！”}
}
{
    \eA{[English source not present in the checked-in Bencraft/Joly files.]}
}
{
}

\Verse{212}
{
    \zA{李纨笑道：“人家不得贵婿，反捱打，我也不 忍得。”}
}
{
    \eA{[English source not present in the checked-in Bencraft/Joly files.]}
}
{
}

\Verse{213}
{
    \zA{众人都笑了。}
}
{
    \eA{[English source not present in the checked-in Bencraft/Joly files.]}
}
{
}

\Verse{214}
{
    \zA{袭人才要掷，只听有人叫门，老婆子忙出去问时，原来是薛姨妈打发人来了接 黛玉的。}
}
{
    \eA{[English source not present in the checked-in Bencraft/Joly files.]}
}
{
}

\Verse{215}
{
    \zA{众人因问：“几更了？”}
}
{
    \eA{[English source not present in the checked-in Bencraft/Joly files.]}
}
{
}

\Verse{216}
{
    \zA{人回：“二更以后了，钟打过十一下了。”}
}
{
    \eA{[English source not present in the checked-in Bencraft/Joly files.]}
}
{
}

\Verse{217}
{
    \zA{宝玉 犹不信，要过表来瞧了一瞧，已是子初一刻十分了，黛玉便起身说：“我可掌不住 了，回去还要吃药呢。”}
}
{
    \eA{[English source not present in the checked-in Bencraft/Joly files.]}
}
{
}

\Verse{218}
{
    \zA{众人说：“也都该散了。”}
}
{
    \eA{[English source not present in the checked-in Bencraft/Joly files.]}
}
{
}

\Verse{219}
{
    \zA{袭人宝玉等还要留着众人，李 纨探春等都说：“夜太深了不像，这已是破格了。”}
}
{
    \eA{[English source not present in the checked-in Bencraft/Joly files.]}
}
{
}

\Verse{220}
{
    \zA{袭人道：“既如此，每位再吃 一杯再走。”}
}
{
    \eA{[English source not present in the checked-in Bencraft/Joly files.]}
}
{
}

\Verse{221}
{
    \zA{说着，晴雯等已都斟满了酒。}
}
{
    \eA{[English source not present in the checked-in Bencraft/Joly files.]}
}
{
}

\Verse{222}
{
    \zA{每人吃了，都命点灯。}
}
{
    \eA{[English source not present in the checked-in Bencraft/Joly files.]}
}
{
}

\Verse{223}
{
    \zA{袭人等齐送过沁 芳亭河那边，方回来。}
}
{
    \eA{[English source not present in the checked-in Bencraft/Joly files.]}
}
{
}

\Verse{224}
{
    \zA{关了门，大家复又行起令来。}
}
{
    \eA{[English source not present in the checked-in Bencraft/Joly files.]}
}
{
}

\Verse{225}
{
    \zA{袭人等又用大钟斟了几钟，用盘子攒了各样果菜 与地下的老妈妈们吃。}
}
{
    \eA{[English source not present in the checked-in Bencraft/Joly files.]}
}
{
}

\Verse{226}
{
    \zA{彼此有了三分酒，便？}
}
{
    \eA{[English source not present in the checked-in Bencraft/Joly files.]}
}
{
}

\Verse{227}
{
    \zA{拳赢唱小曲儿。}
}
{
    \eA{[English source not present in the checked-in Bencraft/Joly files.]}
}
{
}

\Verse{228}
{
    \zA{那天已四更时分，老 妈妈们一面明吃，一面暗偷，酒缸已罄，众人听了，方收拾盥漱睡觉。}
}
{
    \eA{[English source not present in the checked-in Bencraft/Joly files.]}
}
{
}

\Verse{229}
{
    \zA{芳官吃得两 腮胭脂一般，眉梢眼角，添了许多丰韵，身子图不得，便睡在袭人身上，说：“姐 姐，我心跳的很。”}
}
{
    \eA{[English source not present in the checked-in Bencraft/Joly files.]}
}
{
}

\Verse{230}
{
    \zA{袭人笑道：“谁叫你尽力灌呢。”}
}
{
    \eA{[English source not present in the checked-in Bencraft/Joly files.]}
}
{
}

\Verse{231}
{
    \zA{春燕四儿也图不得，早睡了， 晴雯还只管叫。}
}
{
    \eA{[English source not present in the checked-in Bencraft/Joly files.]}
}
{
}

\Verse{232}
{
    \zA{宝玉道：“不用叫了，咱们且胡乱歇一歇。”}
}
{
    \eA{[English source not present in the checked-in Bencraft/Joly files.]}
}
{
}

\Verse{233}
{
    \zA{自己便枕了那红香枕， 身子一歪，就睡着了。}
}
{
    \eA{[English source not present in the checked-in Bencraft/Joly files.]}
}
{
}

\Verse{234}
{
    \zA{袭人见芳官醉的很，恐闹他吐酒，只得轻轻起来，就将芳官 扶在宝玉之侧，由他睡了。}
}
{
    \eA{[English source not present in the checked-in Bencraft/Joly files.]}
}
{
}

\Verse{235}
{
    \zA{自己却在对面榻上倒下。}
}
{
    \eA{[English source not present in the checked-in Bencraft/Joly files.]}
}
{
}

\Verse{236}
{
    \zA{大家黑甜一觉，不知所之。}
}
{
    \eA{[English source not present in the checked-in Bencraft/Joly files.]}
}
{
}

\Verse{237}
{
    \zA{及至天明，袭人睁眼一看，只见天色晶明，忙说： “可迟了！”}
}
{
    \eA{[English source not present in the checked-in Bencraft/Joly files.]}
}
{
}

\Verse{238}
{
    \zA{向对面床上瞧了一瞧，只见芳官头枕着炕沿上，睡犹未醒，连忙起来 叫他。}
}
{
    \eA{[English source not present in the checked-in Bencraft/Joly files.]}
}
{
}

\Verse{239}
{
    \zA{宝玉已翻身醒了。}
}
{
    \eA{[English source not present in the checked-in Bencraft/Joly files.]}
}
{
}

\Verse{240}
{
    \zA{笑道：“可迟了。”}
}
{
    \eA{[English source not present in the checked-in Bencraft/Joly files.]}
}
{
}

\Verse{241}
{
    \zA{因又推芳官起身。}
}
{
    \eA{[English source not present in the checked-in Bencraft/Joly files.]}
}
{
}

\Verse{242}
{
    \zA{那芳官坐起来，犹 发怔揉眼睛。}
}
{
    \eA{[English source not present in the checked-in Bencraft/Joly files.]}
}
{
}

\Verse{243}
{
    \zA{袭人笑道：“不害羞，你喝醉了，怎么也不拣地方儿，乱挺下了？”}
}
{
    \eA{[English source not present in the checked-in Bencraft/Joly files.]}
}
{
}

\Verse{244}
{
    \zA{芳官听了，瞧了瞧，方知是和宝玉同榻，忙羞的笑着下地说：“我怎么――”却说 不出下半句来。}
}
{
    \eA{[English source not present in the checked-in Bencraft/Joly files.]}
}
{
}

\Verse{245}
{
    \zA{宝玉笑道：“我竟也不知道了。}
}
{
    \eA{[English source not present in the checked-in Bencraft/Joly files.]}
}
{
}

\Verse{246}
{
    \zA{若知道，给你脸上抹些墨。”}
}
{
    \eA{[English source not present in the checked-in Bencraft/Joly files.]}
}
{
}

\Verse{247}
{
    \zA{说着， 丫头进来，伺候梳洗。}
}
{
    \eA{[English source not present in the checked-in Bencraft/Joly files.]}
}
{
}

\Verse{248}
{
    \zA{宝玉笑道：“昨日有扰，今日晚上我还席。”}
}
{
    \eA{[English source not present in the checked-in Bencraft/Joly files.]}
}
{
}

\Verse{249}
{
    \zA{袭人笑道：“罢 罢，今日可别闹了，再闹就有人说话了。”}
}
{
    \eA{[English source not present in the checked-in Bencraft/Joly files.]}
}
{
}

\Verse{250}
{
    \zA{宝玉道：“怕什么，不过才两次罢了。}
}
{
    \eA{[English source not present in the checked-in Bencraft/Joly files.]}
}
{
}

\Verse{251}
{
    \zA{咱们也算会吃酒了，一坛子酒怎么就吃光了。}
}
{
    \eA{[English source not present in the checked-in Bencraft/Joly files.]}
}
{
}

\Verse{252}
{
    \zA{――正在有趣儿，偏又没了。”}
}
{
    \eA{[English source not present in the checked-in Bencraft/Joly files.]}
}
{
}

\Verse{253}
{
    \zA{袭人 笑道：“原要这么着才有趣儿，必尽了兴，反无味。}
}
{
    \eA{[English source not present in the checked-in Bencraft/Joly files.]}
}
{
}

\Verse{254}
{
    \zA{昨日都好上来了，晴雯连臊也 忘了，我记得他还唱了一个曲儿。”}
}
{
    \eA{[English source not present in the checked-in Bencraft/Joly files.]}
}
{
}

\Verse{255}
{
    \zA{四儿笑道：“姐姐忘了，连姐姐还唱了一个呢! 在席的谁没唱过？”}
}
{
    \eA{[English source not present in the checked-in Bencraft/Joly files.]}
}
{
}

\Verse{256}
{
    \zA{众人听了，俱红了脸，用两手握着，笑个不住。}
}
{
    \eA{[English source not present in the checked-in Bencraft/Joly files.]}
}
{
}

\Verse{257}
{
    \zA{忽见平儿笑嘻嘻的走来，说：“我亲自来请昨日在席的人，今日我还东，短一 个也使不得。”}
}
{
    \eA{[English source not present in the checked-in Bencraft/Joly files.]}
}
{
}

\Verse{258}
{
    \zA{众人忙让坐吃茶。}
}
{
    \eA{[English source not present in the checked-in Bencraft/Joly files.]}
}
{
}

\Verse{259}
{
    \zA{晴雯笑道：“可惜昨夜没他。”}
}
{
    \eA{[English source not present in the checked-in Bencraft/Joly files.]}
}
{
}

\Verse{260}
{
    \zA{平儿忙问：“你 们夜里做什么来？”}
}
{
    \eA{[English source not present in the checked-in Bencraft/Joly files.]}
}
{
}

\Verse{261}
{
    \zA{袭人便说：“告诉不得你!昨日夜里热闹非常，连往日老太太、 太太带着众人玩，也不及昨儿这一玩：一坛酒我们都鼓捣光了。}
}
{
    \eA{[English source not present in the checked-in Bencraft/Joly files.]}
}
{
}

\Verse{262}
{
    \zA{一个个喝的把臊都 丢了，又都唱起来。}
}
{
    \eA{[English source not present in the checked-in Bencraft/Joly files.]}
}
{
}

\Verse{263}
{
    \zA{四更多天，才横三竖四的打了一个盹儿。”}
}
{
    \eA{[English source not present in the checked-in Bencraft/Joly files.]}
}
{
}

\Verse{264}
{
    \zA{平儿笑道：“好， 白和我要了酒来，也不请我。}
}
{
    \eA{[English source not present in the checked-in Bencraft/Joly files.]}
}
{
}

\Verse{265}
{
    \zA{还说着给我听，气我。”}
}
{
    \eA{[English source not present in the checked-in Bencraft/Joly files.]}
}
{
}

\Verse{266}
{
    \zA{晴雯道：“今儿他还席，必 自来请你，你等着罢。”}
}
{
    \eA{[English source not present in the checked-in Bencraft/Joly files.]}
}
{
}

\Verse{267}
{
    \zA{平儿笑问道：“‘他’是谁？}
}
{
    \eA{[English source not present in the checked-in Bencraft/Joly files.]}
}
{
}

\Verse{268}
{
    \zA{谁是‘他’？”}
}
{
    \eA{[English source not present in the checked-in Bencraft/Joly files.]}
}
{
}

\Verse{269}
{
    \zA{晴雯听了， 把脸飞红了，赶着打，笑说道：“偏你这耳朵尖，听的真！”}
}
{
    \eA{[English source not present in the checked-in Bencraft/Joly files.]}
}
{
}

\Verse{270}
{
    \zA{平儿笑道：“呸!不 害臊的丫头!这会子有事，不和你说。}
}
{
    \eA{[English source not present in the checked-in Bencraft/Joly files.]}
}
{
}

\Verse{271}
{
    \zA{我有事，去了回来再打发人来请。}
}
{
    \eA{[English source not present in the checked-in Bencraft/Joly files.]}
}
{
}

\Verse{272}
{
    \zA{一个不到， 我是打上门来的。”}
}
{
    \eA{[English source not present in the checked-in Bencraft/Joly files.]}
}
{
}

\Verse{273}
{
    \zA{宝玉等忙留他，已经去了。}
}
{
    \eA{[English source not present in the checked-in Bencraft/Joly files.]}
}
{
}

\Verse{274}
{
    \zA{这里宝玉梳洗了，正喝茶，忽然一眼看见砚台底下压着一张纸，因说道：“你 们这么随便混压东西，也不好。”}
}
{
    \eA{[English source not present in the checked-in Bencraft/Joly files.]}
}
{
}

\Verse{275}
{
    \zA{袭人晴雯等忙问：“又怎么了？}
}
{
    \eA{[English source not present in the checked-in Bencraft/Joly files.]}
}
{
}

\Verse{276}
{
    \zA{谁又有了不是了？”}
}
{
    \eA{[English source not present in the checked-in Bencraft/Joly files.]}
}
{
}

\Verse{277}
{
    \zA{宝玉指道：“砚台下是什么？}
}
{
    \eA{[English source not present in the checked-in Bencraft/Joly files.]}
}
{
}

\Verse{278}
{
    \zA{一定又是那位的样子，忘记收的。”}
}
{
    \eA{[English source not present in the checked-in Bencraft/Joly files.]}
}
{
}

\Verse{279}
{
    \zA{晴雯忙启砚拿了 出来，却是一张字帖儿。}
}
{
    \eA{[English source not present in the checked-in Bencraft/Joly files.]}
}
{
}

\Verse{280}
{
    \zA{递给宝玉看时，原来是一张粉红笺纸，上面写着：“槛外 人妙玉恭肃遥叩芳辰。”}
}
{
    \eA{[English source not present in the checked-in Bencraft/Joly files.]}
}
{
}

\Verse{281}
{
    \zA{宝玉看毕，直跳了起来，忙问：“是谁接了来的？}
}
{
    \eA{[English source not present in the checked-in Bencraft/Joly files.]}
}
{
}

\Verse{282}
{
    \zA{也不告 诉！”}
}
{
    \eA{[English source not present in the checked-in Bencraft/Joly files.]}
}
{
}

\Verse{283}
{
    \zA{袭人晴雯等见了这般，不知当是那个要紧的人来的帖子，忙一齐问：“昨儿 是谁接下了一个帖子？”}
}
{
    \eA{[English source not present in the checked-in Bencraft/Joly files.]}
}
{
}

\Verse{284}
{
    \zA{四儿忙跑进来，笑说：“昨儿妙玉并没亲来，只打发个妈 妈送来。}
}
{
    \eA{[English source not present in the checked-in Bencraft/Joly files.]}
}
{
}

\Verse{285}
{
    \zA{我就搁在这里，谁知一顿酒喝的就忘了。”}
}
{
    \eA{[English source not present in the checked-in Bencraft/Joly files.]}
}
{
}

\Verse{286}
{
    \zA{众人听了道：“我当是谁，大 惊小怪，这也不值的。”}
}
{
    \eA{[English source not present in the checked-in Bencraft/Joly files.]}
}
{
}

\Verse{287}
{
    \zA{宝玉忙命：“快拿纸来。”}
}
{
    \eA{[English source not present in the checked-in Bencraft/Joly files.]}
}
{
}

\Verse{288}
{
    \zA{当下拿了纸，研了墨，看他下 着“槛外人”三字，自己竟不知回帖上回个什么字样才相敌，只管提笔出神，半天 仍没主意。}
}
{
    \eA{[English source not present in the checked-in Bencraft/Joly files.]}
}
{
}

\Verse{289}
{
    \zA{因又想：“要问宝钗去，他必又批评怪诞，不如问黛玉去。”}
}
{
    \eA{[English source not present in the checked-in Bencraft/Joly files.]}
}
{
}

\Verse{290}
{
    \zA{想罢，袖 了帖儿，径来寻黛玉。}
}
{
    \eA{[English source not present in the checked-in Bencraft/Joly files.]}
}
{
}

\Verse{291}
{
    \zA{刚过了沁芳亭，忽见岫烟颤颤巍巍的迎面走来。}
}
{
    \eA{[English source not present in the checked-in Bencraft/Joly files.]}
}
{
}

\Verse{292}
{
    \zA{宝玉忙问：“姐姐那里去？”}
}
{
    \eA{[English source not present in the checked-in Bencraft/Joly files.]}
}
{
}

\Verse{293}
{
    \zA{岫烟笑道：“我找妙玉说话。”}
}
{
    \eA{[English source not present in the checked-in Bencraft/Joly files.]}
}
{
}

\Verse{294}
{
    \zA{宝玉听了，诧异说道：“他为人孤癖，不合时宜， 万人不入他的目。}
}
{
    \eA{[English source not present in the checked-in Bencraft/Joly files.]}
}
{
}

\Verse{295}
{
    \zA{原来他推重姐姐，竟知姐姐不是我们一流俗人。”}
}
{
    \eA{[English source not present in the checked-in Bencraft/Joly files.]}
}
{
}

\Verse{296}
{
    \zA{岫烟笑道：“他 也未必真心重我，但我和他做过十年的邻居，只一墙之隔。}
}
{
    \eA{[English source not present in the checked-in Bencraft/Joly files.]}
}
{
}

\Verse{297}
{
    \zA{他在蟠香寺修炼，我家 原来寒素，赁房居就，赁了他庙里的房子住了十年。}
}
{
    \eA{[English source not present in the checked-in Bencraft/Joly files.]}
}
{
}

\Verse{298}
{
    \zA{无事到他庙里去作伴，我所认 得的字，都是承他所授：我和他又是贫贱之交，又有半师之分。}
}
{
    \eA{[English source not present in the checked-in Bencraft/Joly files.]}
}
{
}

\Verse{299}
{
    \zA{因我们投亲去了， 闻得他因不合时宜，权势不容，竟投到这里来。}
}
{
    \eA{[English source not present in the checked-in Bencraft/Joly files.]}
}
{
}

\Verse{300}
{
    \zA{如今又两缘凑合，我们得遇，旧情 竟未改易，承他青目，更胜当日。”}
}
{
    \eA{[English source not present in the checked-in Bencraft/Joly files.]}
}
{
}

\Verse{301}
{
    \zA{宝玉听了，恍如听了焦雷一般，喜得笑道：“怪 道姐姐举止言谈，超然如野鹤闲云，原本有来历。}
}
{
    \eA{[English source not present in the checked-in Bencraft/Joly files.]}
}
{
}

\Verse{302}
{
    \zA{我正因他的一件事为难，要请教 别人去。}
}
{
    \eA{[English source not present in the checked-in Bencraft/Joly files.]}
}
{
}

\Verse{303}
{
    \zA{如今遇见姐姐，真是天缘凑合，求姐姐指教。”}
}
{
    \eA{[English source not present in the checked-in Bencraft/Joly files.]}
}
{
}

\Verse{304}
{
    \zA{说着便将拜帖取给岫烟看。}
}
{
    \eA{[English source not present in the checked-in Bencraft/Joly files.]}
}
{
}

\Verse{305}
{
    \zA{岫烟笑道：“他这脾气竟不能改，竟是生成这等放诞诡僻了。}
}
{
    \eA{[English source not present in the checked-in Bencraft/Joly files.]}
}
{
}

\Verse{306}
{
    \zA{从来没见拜帖上下别 号的，这可是俗语说的‘僧不僧，俗不俗，女不女，男不男’，成个什么理数。”}
}
{
    \eA{[English source not present in the checked-in Bencraft/Joly files.]}
}
{
}

\Verse{307}
{
    \zA{宝玉听说，忙笑道：“姐姐不知道，他原不在这些人中里，他原是世人意外之人。}
}
{
    \eA{[English source not present in the checked-in Bencraft/Joly files.]}
}
{
}

\Verse{308}
{
    \zA{因取了我是个些微有知识的，方给我这帖子。}
}
{
    \eA{[English source not present in the checked-in Bencraft/Joly files.]}
}
{
}

\Verse{309}
{
    \zA{我因不知回什么字样才好，竟没了主 意，正要去问林妹妹，可巧遇见了姐姐。”}
}
{
    \eA{[English source not present in the checked-in Bencraft/Joly files.]}
}
{
}

\Verse{310}
{
    \zA{岫烟听了宝玉这话，且只管用眼上下细细打量了半日，方笑道：“怪道俗语说 的，‘闻名不如见面’，又怪不的妙玉竟下这帖子给你，又怪不的上年竟给你那些 梅花。}
}
{
    \eA{[English source not present in the checked-in Bencraft/Joly files.]}
}
{
}

\Verse{311}
{
    \zA{既连他这样，少不得我告诉你原故。}
}
{
    \eA{[English source not present in the checked-in Bencraft/Joly files.]}
}
{
}

\Verse{312}
{
    \zA{他常说古人中自汉、晋、五代、唐、宋 以来，皆无好诗，只有两句好，说道：‘纵有千年铁门槛，终须一个土馒头。’}
}
{
    \eA{[English source not present in the checked-in Bencraft/Joly files.]}
}
{
}

\Verse{313}
{
    \zA{所 以他自称‘槛外之人’。}
}
{
    \eA{[English source not present in the checked-in Bencraft/Joly files.]}
}
{
}

\Verse{314}
{
    \zA{又常赞：‘文是庄子的好。’}
}
{
    \eA{[English source not present in the checked-in Bencraft/Joly files.]}
}
{
}

\Verse{315}
{
    \zA{故又或称为‘畸人’。}
}
{
    \eA{[English source not present in the checked-in Bencraft/Joly files.]}
}
{
}

\Verse{316}
{
    \zA{他若 帖子上是自称‘畸人’的，你就还他个‘世人’。}
}
{
    \eA{[English source not present in the checked-in Bencraft/Joly files.]}
}
{
}

\Verse{317}
{
    \zA{‘畸人’者，他自称是畸零之人， 你谦自己乃世人扰扰之人，他便喜了。}
}
{
    \eA{[English source not present in the checked-in Bencraft/Joly files.]}
}
{
}

\Verse{318}
{
    \zA{如今他自称‘槛外之人’，是自谓蹈于铁槛 之外了，故你如今只下‘槛内人’，便合了他的心了。”}
}
{
    \eA{[English source not present in the checked-in Bencraft/Joly files.]}
}
{
}

\Verse{319}
{
    \zA{宝玉听了，如醍醐灌顶， “嗳哟”了一声，方笑道：“怪道我们家庙说是铁槛寺呢，原来有这一说。}
}
{
    \eA{[English source not present in the checked-in Bencraft/Joly files.]}
}
{
}

\Verse{320}
{
    \zA{姐姐就 请，让我去写回帖。”}
}
{
    \eA{[English source not present in the checked-in Bencraft/Joly files.]}
}
{
}

\Verse{321}
{
    \zA{岫烟听了，便自往栊翠庵来。}
}
{
    \eA{[English source not present in the checked-in Bencraft/Joly files.]}
}
{
}

\Verse{322}
{
    \zA{宝玉回房，写了帖子，上面只 写“槛内人宝玉熏沐谨拜”几字。}
}
{
    \eA{[English source not present in the checked-in Bencraft/Joly files.]}
}
{
}

\Verse{323}
{
    \zA{亲自拿了到栊翠庵，只隔门缝儿投进去，便回来 了。}
}
{
    \eA{[English source not present in the checked-in Bencraft/Joly files.]}
}
{
}

\Verse{324}
{
    \zA{因饭后平儿还席，说红香圃太热，便在榆荫堂中摆了几席新酒佳肴。}
}
{
    \eA{[English source not present in the checked-in Bencraft/Joly files.]}
}
{
}

\Verse{325}
{
    \zA{可喜尤氏 又带了佩凤偕鸾二妾过来游玩。}
}
{
    \eA{[English source not present in the checked-in Bencraft/Joly files.]}
}
{
}

\Verse{326}
{
    \zA{这二妾亦是青年娇憨女子，不常过来的，今既入了 这园，再遇见湘云、香菱、芳、蕊一干女子，所谓“方以类聚，物以群分”二语不
        错，只见他们说笑不了，也不管尤氏在那里，只凭丫鬟们去服役，且同众人一一的 游玩。}
}
{
    \eA{[English source not present in the checked-in Bencraft/Joly files.]}
}
{
}

\Verse{327}
{
    \zA{闲言少述，且说当下众人都在榆荫堂中，以酒为名，大家玩笑，命女先儿击鼓。}
}
{
    \eA{[English source not present in the checked-in Bencraft/Joly files.]}
}
{
}

\Verse{328}
{
    \zA{平儿采了一枝芍药，大家约二十来人，传花为令，热闹了一回。}
}
{
    \eA{[English source not present in the checked-in Bencraft/Joly files.]}
}
{
}

\Verse{329}
{
    \zA{因人回说：“甄家 有两个女人送东西来了。”}
}
{
    \eA{[English source not present in the checked-in Bencraft/Joly files.]}
}
{
}

\Verse{330}
{
    \zA{探春和李纨尤氏三人出去议事厅相见。}
}
{
    \eA{[English source not present in the checked-in Bencraft/Joly files.]}
}
{
}

\Verse{331}
{
    \zA{这里众人且出来 散一散。}
}
{
    \eA{[English source not present in the checked-in Bencraft/Joly files.]}
}
{
}

\Verse{332}
{
    \zA{佩凤偕鸾两个去打秋千玩耍，宝玉便说：“你两个上去，让我送。”}
}
{
    \eA{[English source not present in the checked-in Bencraft/Joly files.]}
}
{
}

\Verse{333}
{
    \zA{慌的 佩凤说：“罢了，别替我们闹乱子！”}
}
{
    \eA{[English source not present in the checked-in Bencraft/Joly files.]}
}
{
}

\Verse{334}
{
    \zA{忽见东府里几个人，慌慌张张跑来，说：“老爷殡天了！”}
}
{
    \eA{[English source not present in the checked-in Bencraft/Joly files.]}
}
{
}

\Verse{335}
{
    \zA{众人听了，吓了一 大跳，忙都说：“好好的并无疾病，怎么就没了？”}
}
{
    \eA{[English source not present in the checked-in Bencraft/Joly files.]}
}
{
}

\Verse{336}
{
    \zA{家人说：“老爷天天修炼，定 是功成圆满，升仙去了。”}
}
{
    \eA{[English source not present in the checked-in Bencraft/Joly files.]}
}
{
}

\Verse{337}
{
    \zA{尤氏一闻此言，又见贾珍父子并贾琏等皆不在家，一时 竟没个着己的男子来，未免忙了。}
}
{
    \eA{[English source not present in the checked-in Bencraft/Joly files.]}
}
{
}

\Verse{338}
{
    \zA{只得忙卸了妆饰，命人先到玄真观将所有的道士 都锁了起来，等大爷来家审问；一面忙忙坐车，带了赖升一干老人媳妇出城。}
}
{
    \eA{[English source not present in the checked-in Bencraft/Joly files.]}
}
{
}

\Verse{339}
{
    \zA{又请 大夫看视，到底系何病症。}
}
{
    \eA{[English source not present in the checked-in Bencraft/Joly files.]}
}
{
}

\Verse{340}
{
    \zA{大夫们见人已死，何处诊脉来？}
}
{
    \eA{[English source not present in the checked-in Bencraft/Joly files.]}
}
{
}

\Verse{341}
{
    \zA{素知贾敬导气之术，总 属虚诞，更至参星礼斗，守庚申，服灵砂等，妄作虚为，过于劳神费力，反因此伤
        了性命的，如今虽死，腹中坚硬似铁，面皮嘴唇，烧的紫绛皱裂。}
}
{
    \eA{[English source not present in the checked-in Bencraft/Joly files.]}
}
{
}

\Verse{342}
{
    \zA{便向媳妇回说： “系道教中吞金服砂，烧胀而殁。”}
}
{
    \eA{[English source not present in the checked-in Bencraft/Joly files.]}
}
{
}

\Verse{343}
{
    \zA{众道士慌的回道：“原是秘制的丹砂吃坏了事， 小道们也曾劝说：‘功夫未到，且服不得。’}
}
{
    \eA{[English source not present in the checked-in Bencraft/Joly files.]}
}
{
}

\Verse{344}
{
    \zA{不承望老爷于今夜守庚申时，悄悄的 服了下去，便升仙去了。}
}
{
    \eA{[English source not present in the checked-in Bencraft/Joly files.]}
}
{
}

\Verse{345}
{
    \zA{这是虔心得道，已出苦海，脱去皮囊了。”}
}
{
    \eA{[English source not present in the checked-in Bencraft/Joly files.]}
}
{
}

\Verse{346}
{
    \zA{尤氏也不便听， 只命锁着，等贾珍来发放，且命人飞马报信。}
}
{
    \eA{[English source not present in the checked-in Bencraft/Joly files.]}
}
{
}

\Verse{347}
{
    \zA{一面看视里面窄狭，不能停放，横竖 也不能进城的，忙装裹好了，用软轿抬至铁槛寺来停放。}
}
{
    \eA{[English source not present in the checked-in Bencraft/Joly files.]}
}
{
}

\Verse{348}
{
    \zA{掐指算来，至早也得半月 的工夫贾珍方能来到，目今天气炎热，实不能相待，遂自行主持，命天文生择了日 期入殓。}
}
{
    \eA{[English source not present in the checked-in Bencraft/Joly files.]}
}
{
}

\Verse{349}
{
    \zA{寿木早年已经备下，寄在此庙的，甚是便宜。}
}
{
    \eA{[English source not present in the checked-in Bencraft/Joly files.]}
}
{
}

\Verse{350}
{
    \zA{三日后，便破孝开吊，一面 且做起道场来。}
}
{
    \eA{[English source not present in the checked-in Bencraft/Joly files.]}
}
{
}

\Verse{351}
{
    \zA{因那边荣府里凤姐儿出不来，李纨又照顾姐妹，宝玉不识事体，只 得将外头事务，暂托了几个家里二等管事的。}
}
{
    \eA{[English source not present in the checked-in Bencraft/Joly files.]}
}
{
}

\Verse{352}
{
    \zA{贾？}
}
{
    \eA{[English source not present in the checked-in Bencraft/Joly files.]}
}
{
}

\Verse{353}
{
    \zA{、贾？}
}
{
    \eA{[English source not present in the checked-in Bencraft/Joly files.]}
}
{
}

\Verse{354}
{
    \zA{、贾珩、贾璎、贾菖、贾 菱等各有执事。}
}
{
    \eA{[English source not present in the checked-in Bencraft/Joly files.]}
}
{
}

\Verse{355}
{
    \zA{尤氏不能回家，便将他继母接来，在宁府看家。}
}
{
    \eA{[English source not present in the checked-in Bencraft/Joly files.]}
}
{
}

\Verse{356}
{
    \zA{这继母只得将两个 未出嫁的女儿带来，一并住着，才放心。}
}
{
    \eA{[English source not present in the checked-in Bencraft/Joly files.]}
}
{
}

\Verse{357}
{
    \zA{且说贾珍闻了此信，急忙告假，――并贾蓉是有职人员。}
}
{
    \eA{[English source not present in the checked-in Bencraft/Joly files.]}
}
{
}

\Verse{358}
{
    \zA{礼部见当今隆敦孝弟， 不敢自专，具本请旨。}
}
{
    \eA{[English source not present in the checked-in Bencraft/Joly files.]}
}
{
}

\Verse{359}
{
    \zA{原来天子极是仁孝过天的，且更隆重功臣之裔，一见此本， 便诏问贾敬何职。}
}
{
    \eA{[English source not present in the checked-in Bencraft/Joly files.]}
}
{
}

\Verse{360}
{
    \zA{礼部代奏：“系进士出身，祖职已荫其子贾珍。}
}
{
    \eA{[English source not present in the checked-in Bencraft/Joly files.]}
}
{
}

\Verse{361}
{
    \zA{贾敬因年迈多疾， 常养静于都城之外玄真观，今因疾殁于观中。}
}
{
    \eA{[English source not present in the checked-in Bencraft/Joly files.]}
}
{
}

\Verse{362}
{
    \zA{其子珍，其孙蓉，现因国丧，随驾在 此，故乞假归殓。”}
}
{
    \eA{[English source not present in the checked-in Bencraft/Joly files.]}
}
{
}

\Verse{363}
{
    \zA{天子听了，忙下额外恩旨曰：“贾敬虽无功于国，念彼祖父之 忠，追赐五品之职。}
}
{
    \eA{[English source not present in the checked-in Bencraft/Joly files.]}
}
{
}

\Verse{364}
{
    \zA{令其子孙扶柩由北下门入都，恩赐私第殡殓，任子孙尽丧，礼 毕扶柩回籍。}
}
{
    \eA{[English source not present in the checked-in Bencraft/Joly files.]}
}
{
}

\Verse{365}
{
    \zA{外着光禄寺按上例赐祭，朝中由王公以下，准其祭吊。}
}
{
    \eA{[English source not present in the checked-in Bencraft/Joly files.]}
}
{
}

\Verse{366}
{
    \zA{钦此。”}
}
{
    \eA{[English source not present in the checked-in Bencraft/Joly files.]}
}
{
}

\Verse{367}
{
    \zA{此旨 一下，不但贾府里人谢恩，连朝中所有大臣，皆嵩呼称颂不绝。}
}
{
    \eA{[English source not present in the checked-in Bencraft/Joly files.]}
}
{
}

\Verse{368}
{
    \zA{贾珍父子星夜驰回。}
}
{
    \eA{[English source not present in the checked-in Bencraft/Joly files.]}
}
{
}

\Verse{369}
{
    \zA{半路中又见贾？}
}
{
    \eA{[English source not present in the checked-in Bencraft/Joly files.]}
}
{
}

\Verse{370}
{
    \zA{贾？}
}
{
    \eA{[English source not present in the checked-in Bencraft/Joly files.]}
}
{
}

\Verse{371}
{
    \zA{二人，领家丁飞骑而来，看见贾珍， 一齐滚鞍下马请安。}
}
{
    \eA{[English source not present in the checked-in Bencraft/Joly files.]}
}
{
}

\Verse{372}
{
    \zA{贾珍忙问：“做什么？”}
}
{
    \eA{[English source not present in the checked-in Bencraft/Joly files.]}
}
{
}

\Verse{373}
{
    \zA{贾？}
}
{
    \eA{[English source not present in the checked-in Bencraft/Joly files.]}
}
{
}

\Verse{374}
{
    \zA{回说：“嫂子恐哥哥和侄儿来了， 老太太路上无人，叫我们两个来护送老太太的。”}
}
{
    \eA{[English source not present in the checked-in Bencraft/Joly files.]}
}
{
}

\Verse{375}
{
    \zA{贾珍听了，赞声不绝。}
}
{
    \eA{[English source not present in the checked-in Bencraft/Joly files.]}
}
{
}

\Verse{376}
{
    \zA{又问：“家 中如何料理？”}
}
{
    \eA{[English source not present in the checked-in Bencraft/Joly files.]}
}
{
}

\Verse{377}
{
    \zA{贾？}
}
{
    \eA{[English source not present in the checked-in Bencraft/Joly files.]}
}
{
}

\Verse{378}
{
    \zA{等便将如何拿了道士，如何挪至家庙，怕家内无人，接了亲家 母和两个姨奶奶在上房住着。}
}
{
    \eA{[English source not present in the checked-in Bencraft/Joly files.]}
}
{
}

\Verse{379}
{
    \zA{贾蓉当下也下了马，听见两个姨娘来了，喜的笑容满 面。}
}
{
    \eA{[English source not present in the checked-in Bencraft/Joly files.]}
}
{
}

\Verse{380}
{
    \zA{贾珍忙说了几声“妥当”，加鞭便走。}
}
{
    \eA{[English source not present in the checked-in Bencraft/Joly files.]}
}
{
}

\Verse{381}
{
    \zA{店也不投，连夜换马飞驰。}
}
{
    \eA{[English source not present in the checked-in Bencraft/Joly files.]}
}
{
}

\Verse{382}
{
    \zA{一日到了都 门，先奔入铁槛寺，那天已是四更天气。}
}
{
    \eA{[English source not present in the checked-in Bencraft/Joly files.]}
}
{
}

\Verse{383}
{
    \zA{坐更的闻知，忙喝起众人来。}
}
{
    \eA{[English source not present in the checked-in Bencraft/Joly files.]}
}
{
}

\Verse{384}
{
    \zA{贾珍下了马， 和贾蓉放声大哭，从大门外便跪爬起来，至棺前稽颡泣血，直哭到天亮，喉咙都哭 哑了方住。}
}
{
    \eA{[English source not present in the checked-in Bencraft/Joly files.]}
}
{
}

\Verse{385}
{
    \zA{尤氏等都一齐见过，贾珍父子忙按礼换了凶服，在棺前俯伏。}
}
{
    \eA{[English source not present in the checked-in Bencraft/Joly files.]}
}
{
}

\Verse{386}
{
    \zA{无奈自要 理事，竟不能目不视物、耳不闻声，少不得减了些悲戚，好指挥众人。}
}
{
    \eA{[English source not present in the checked-in Bencraft/Joly files.]}
}
{
}

\Verse{387}
{
    \zA{因将恩旨备 述给众亲友听了，一面先打发贾蓉回家来，料理停灵之事。}
}
{
    \eA{[English source not present in the checked-in Bencraft/Joly files.]}
}
{
}

\Verse{388}
{
    \zA{贾蓉巴不得一声儿，便先骑马跑来。}
}
{
    \eA{[English source not present in the checked-in Bencraft/Joly files.]}
}
{
}

\Verse{389}
{
    \zA{到家，忙命前厅收桌椅，下？}
}
{
    \eA{[English source not present in the checked-in Bencraft/Joly files.]}
}
{
}

\Verse{390}
{
    \zA{扇，挂孝幔 子，门前起鼓手棚、牌楼等事。}
}
{
    \eA{[English source not present in the checked-in Bencraft/Joly files.]}
}
{
}

\Verse{391}
{
    \zA{又忙着进来看外祖母、两个姨娘。}
}
{
    \eA{[English source not present in the checked-in Bencraft/Joly files.]}
}
{
}

\Verse{392}
{
    \zA{原来尤老安人年 高喜睡，常常歪着；他二姨娘三姨娘都和丫头们做活计，见他来了，都道烦恼。}
}
{
    \eA{[English source not present in the checked-in Bencraft/Joly files.]}
}
{
}

\Verse{393}
{
    \zA{贾 蓉且嘻嘻的望他二姨娘笑说：“二姨娘，你又来了？}
}
{
    \eA{[English source not present in the checked-in Bencraft/Joly files.]}
}
{
}

\Verse{394}
{
    \zA{我父亲正想你呢。”}
}
{
    \eA{[English source not present in the checked-in Bencraft/Joly files.]}
}
{
}

\Verse{395}
{
    \zA{二姨娘红 了脸，骂道：“好蓉小子!我过两日不骂你几句，你就过不得了，越发连个体统都 没了。}
}
{
    \eA{[English source not present in the checked-in Bencraft/Joly files.]}
}
{
}

\Verse{396}
{
    \zA{还亏你是大家公子哥儿，每日念书学礼的，越发连那小家子的也跟不上。”}
}
{
    \eA{[English source not present in the checked-in Bencraft/Joly files.]}
}
{
}

\Verse{397}
{
    \zA{说着顺手拿起一个熨斗来，兜头就打，吓得贾蓉抱着头，滚到怀里告饶。}
}
{
    \eA{[English source not present in the checked-in Bencraft/Joly files.]}
}
{
}

\Verse{398}
{
    \zA{尤三姐便 转过脸去，说道：“等姐姐来家再告诉他。”}
}
{
    \eA{[English source not present in the checked-in Bencraft/Joly files.]}
}
{
}

\Verse{399}
{
    \zA{贾蓉忙笑着跪在炕上求饶，因又和他 二姨娘抢砂仁吃。}
}
{
    \eA{[English source not present in the checked-in Bencraft/Joly files.]}
}
{
}

\Verse{400}
{
    \zA{那二姐儿嚼了一嘴渣子，吐了他一脸，贾蓉用舌头都舔着吃了。}
}
{
    \eA{[English source not present in the checked-in Bencraft/Joly files.]}
}
{
}

\Verse{401}
{
    \zA{众丫头看不过，都笑说：“热孝在身上，老娘才睡了觉。}
}
{
    \eA{[English source not present in the checked-in Bencraft/Joly files.]}
}
{
}

\Verse{402}
{
    \zA{他两个虽小，到底是姨娘 家。}
}
{
    \eA{[English source not present in the checked-in Bencraft/Joly files.]}
}
{
}

\Verse{403}
{
    \zA{你太眼里没有奶奶了，回来告诉爷，你吃不了兜着走。”}
}
{
    \eA{[English source not present in the checked-in Bencraft/Joly files.]}
}
{
}

\Verse{404}
{
    \zA{贾蓉撇下他姨娘，便 抱着那丫头亲嘴，说：“我的心肝，你说得是。}
}
{
    \eA{[English source not present in the checked-in Bencraft/Joly files.]}
}
{
}

\Verse{405}
{
    \zA{咱们馋他们两个。”}
}
{
    \eA{[English source not present in the checked-in Bencraft/Joly files.]}
}
{
}

\Verse{406}
{
    \zA{丫头们忙推他， 恨的骂：“短命鬼!你一般有老婆丫头，只和我们闹。}
}
{
    \eA{[English source not present in the checked-in Bencraft/Joly files.]}
}
{
}

\Verse{407}
{
    \zA{知道的说是玩，不知道的人， 再遇见那样脏心烂肺的、爱多管闲事嚼舌头的人，吵嚷到那府里，背地嚼舌，说咱 们这边混账。”}
}
{
    \eA{[English source not present in the checked-in Bencraft/Joly files.]}
}
{
}

\Verse{408}
{
    \zA{贾蓉笑道：“各门另户，谁管谁的事？}
}
{
    \eA{[English source not present in the checked-in Bencraft/Joly files.]}
}
{
}

\Verse{409}
{
    \zA{都够使的了。}
}
{
    \eA{[English source not present in the checked-in Bencraft/Joly files.]}
}
{
}

\Verse{410}
{
    \zA{从古至今，连 汉朝和唐朝，人还说‘脏唐臭汉’，何况咱们这宗人家!谁家没风流事？}
}
{
    \eA{[English source not present in the checked-in Bencraft/Joly files.]}
}
{
}

\Verse{411}
{
    \zA{别叫我说出 来。}
}
{
    \eA{[English source not present in the checked-in Bencraft/Joly files.]}
}
{
}

\Verse{412}
{
    \zA{连那边大老爷这么利害，琏二叔还和那小姨娘不干净呢。}
}
{
    \eA{[English source not present in the checked-in Bencraft/Joly files.]}
}
{
}

\Verse{413}
{
    \zA{凤婶子那样刚强，瑞 大叔还想他的账：那一件瞒了我？”}
}
{
    \eA{[English source not present in the checked-in Bencraft/Joly files.]}
}
{
}

\Verse{414}
{
    \zA{贾蓉只管信口开河，胡言乱道。}
}
{
    \eA{[English source not present in the checked-in Bencraft/Joly files.]}
}
{
}

\Verse{415}
{
    \zA{三姐儿？}
}
{
    \eA{[English source not present in the checked-in Bencraft/Joly files.]}
}
{
}

\Verse{416}
{
    \zA{了脸，早下炕进里间屋里，叫醒尤老 娘。}
}
{
    \eA{[English source not present in the checked-in Bencraft/Joly files.]}
}
{
}

\Verse{417}
{
    \zA{这里贾蓉见他老娘醒了，忙去请安问好。}
}
{
    \eA{[English source not present in the checked-in Bencraft/Joly files.]}
}
{
}

\Verse{418}
{
    \zA{又说：“老祖宗劳心，又难为两位姨 娘受委屈，我们爷儿们感激不尽。}
}
{
    \eA{[English source not present in the checked-in Bencraft/Joly files.]}
}
{
}

\Verse{419}
{
    \zA{惟有等事完了，我们合家大小登门磕头去。”}
}
{
    \eA{[English source not present in the checked-in Bencraft/Joly files.]}
}
{
}

\Verse{420}
{
    \zA{尤 老安人点头道：“我的儿，倒是你会说话。}
}
{
    \eA{[English source not present in the checked-in Bencraft/Joly files.]}
}
{
}

\Verse{421}
{
    \zA{亲戚们原是该的。”}
}
{
    \eA{[English source not present in the checked-in Bencraft/Joly files.]}
}
{
}

\Verse{422}
{
    \zA{又问：“你父亲好？}
}
{
    \eA{[English source not present in the checked-in Bencraft/Joly files.]}
}
{
}

\Verse{423}
{
    \zA{几时得了信赶到的？”}
}
{
    \eA{[English source not present in the checked-in Bencraft/Joly files.]}
}
{
}

\Verse{424}
{
    \zA{贾蓉笑道：“刚才赶到的，先打发我瞧你老人家来了，好歹 求你老人家事完了再去。”}
}
{
    \eA{[English source not present in the checked-in Bencraft/Joly files.]}
}
{
}

\Verse{425}
{
    \zA{说着，又和他二姨娘挤眼儿。}
}
{
    \eA{[English source not present in the checked-in Bencraft/Joly files.]}
}
{
}

\Verse{426}
{
    \zA{二姐便悄悄咬牙骂道：“很 会嚼舌根的猴儿崽子!留下我们，给你爹做妈不成？”}
}
{
    \eA{[English source not present in the checked-in Bencraft/Joly files.]}
}
{
}

\Verse{427}
{
    \zA{贾蓉又和尤老娘道：“放心 罢，我父亲每日为两位姨娘操心，要寻两个有根基的富贵人家，又年轻又俏皮两位 姨娘父亲，好聘嫁这二位姨娘。}
}
{
    \eA{[English source not present in the checked-in Bencraft/Joly files.]}
}
{
}

\Verse{428}
{
    \zA{这几年总没拣着，可巧前儿路上才相准了一个。”}
}
{
    \eA{[English source not present in the checked-in Bencraft/Joly files.]}
}
{
}

\Verse{429}
{
    \zA{尤老娘只当是真话，忙问：“是谁家的？”}
}
{
    \eA{[English source not present in the checked-in Bencraft/Joly files.]}
}
{
}

\Verse{430}
{
    \zA{二姐丢了活计，一头笑，一头赶着打， 说：“妈妈，别信这混账孩子的话。”}
}
{
    \eA{[English source not present in the checked-in Bencraft/Joly files.]}
}
{
}

\Verse{431}
{
    \zA{三姐儿道：“蓉儿，你说是说，别只管嘴里 这么不清不浑的！”}
}
{
    \eA{[English source not present in the checked-in Bencraft/Joly files.]}
}
{
}

\Verse{432}
{
    \zA{说着，人来回话，说：“事已完了，请哥儿出去看了，回爷的 话去呢。”}
}
{
    \eA{[English source not present in the checked-in Bencraft/Joly files.]}
}
{
}

\Verse{433}
{
    \zA{那贾蓉方笑嘻嘻的出来。}
}
{
    \eA{[English source not present in the checked-in Bencraft/Joly files.]}
}
{
}

\Verse{434}
{
    \zA{不知如何，下回分解。}
}
{
    \eA{[English source not present in the checked-in Bencraft/Joly files.]}
}
{
}

\Verse{435}
{
    \zA{(本章完)}
}
{
    \eA{[English source not present in the checked-in Bencraft/Joly files.]}
}
{
}
